% Préambule commun des schémas TikZ — charte Liouaï
% Chaque schéma commence par % Préambule commun des schémas TikZ — charte Liouaï
% Chaque schéma commence par % Préambule commun des schémas TikZ — charte Liouaï
% Chaque schéma commence par \input{_preambule} puis \begin{document}.
\documentclass[tikz,border=6pt]{standalone}
\usepackage[T1]{fontenc}
\usepackage[utf8]{inputenc}
\usepackage[scaled=0.95]{helvet}
\renewcommand{\familydefault}{\sfdefault}
\usepackage{amsmath,amssymb}
\usepackage[european]{circuitikz}
\usetikzlibrary{arrows.meta,positioning,calc,angles,quotes,decorations.pathreplacing,backgrounds}

% Couleurs du logo Liouaï Digital Education
\definecolor{navy}{HTML}{243743}
\definecolor{vert}{HTML}{28B78D}
\definecolor{vertfonce}{HTML}{167A5C}
\definecolor{vertclair}{HTML}{EAF7F2}
\definecolor{gris}{HTML}{A1A9AE}
\definecolor{gristexte}{HTML}{5C6870}
\definecolor{fond}{HTML}{F3F5F6}
\definecolor{jaune}{HTML}{E5C800}
\definecolor{jauneclair}{HTML}{FFFBD6}

\tikzset{
  every picture/.style={line cap=round, line join=round, color=navy, text=navy},
  fleche/.style={-{Stealth[length=2.6mm,width=2mm]}, line width=1.1pt},
  vecteur/.style={fleche, line width=1.4pt},
  titre/.style={font=\bfseries\small, text=navy},
  note/.style={font=\footnotesize, text=gristexte},
  boite/.style={draw=navy, line width=0.9pt, rounded corners=5pt, fill=fond, align=center, inner sep=6pt},
}
 puis \begin{document}.
\documentclass[tikz,border=6pt]{standalone}
\usepackage[T1]{fontenc}
\usepackage[utf8]{inputenc}
\usepackage[scaled=0.95]{helvet}
\renewcommand{\familydefault}{\sfdefault}
\usepackage{amsmath,amssymb}
\usepackage[european]{circuitikz}
\usetikzlibrary{arrows.meta,positioning,calc,angles,quotes,decorations.pathreplacing,backgrounds}

% Couleurs du logo Liouaï Digital Education
\definecolor{navy}{HTML}{243743}
\definecolor{vert}{HTML}{28B78D}
\definecolor{vertfonce}{HTML}{167A5C}
\definecolor{vertclair}{HTML}{EAF7F2}
\definecolor{gris}{HTML}{A1A9AE}
\definecolor{gristexte}{HTML}{5C6870}
\definecolor{fond}{HTML}{F3F5F6}
\definecolor{jaune}{HTML}{E5C800}
\definecolor{jauneclair}{HTML}{FFFBD6}

\tikzset{
  every picture/.style={line cap=round, line join=round, color=navy, text=navy},
  fleche/.style={-{Stealth[length=2.6mm,width=2mm]}, line width=1.1pt},
  vecteur/.style={fleche, line width=1.4pt},
  titre/.style={font=\bfseries\small, text=navy},
  note/.style={font=\footnotesize, text=gristexte},
  boite/.style={draw=navy, line width=0.9pt, rounded corners=5pt, fill=fond, align=center, inner sep=6pt},
}
 puis \begin{document}.
\documentclass[tikz,border=6pt]{standalone}
\usepackage[T1]{fontenc}
\usepackage[utf8]{inputenc}
\usepackage[scaled=0.95]{helvet}
\renewcommand{\familydefault}{\sfdefault}
\usepackage{amsmath,amssymb}
\usepackage[european]{circuitikz}
\usetikzlibrary{arrows.meta,positioning,calc,angles,quotes,decorations.pathreplacing,backgrounds}

% Couleurs du logo Liouaï Digital Education
\definecolor{navy}{HTML}{243743}
\definecolor{vert}{HTML}{28B78D}
\definecolor{vertfonce}{HTML}{167A5C}
\definecolor{vertclair}{HTML}{EAF7F2}
\definecolor{gris}{HTML}{A1A9AE}
\definecolor{gristexte}{HTML}{5C6870}
\definecolor{fond}{HTML}{F3F5F6}
\definecolor{jaune}{HTML}{E5C800}
\definecolor{jauneclair}{HTML}{FFFBD6}

\tikzset{
  every picture/.style={line cap=round, line join=round, color=navy, text=navy},
  fleche/.style={-{Stealth[length=2.6mm,width=2mm]}, line width=1.1pt},
  vecteur/.style={fleche, line width=1.4pt},
  titre/.style={font=\bfseries\small, text=navy},
  note/.style={font=\footnotesize, text=gristexte},
  boite/.style={draw=navy, line width=0.9pt, rounded corners=5pt, fill=fond, align=center, inner sep=6pt},
}

\begin{document}
\begin{tikzpicture}[x={(1.6cm,-0.45cm)}, y={(1.3cm,0.75cm)}, z={(0cm,1.75cm)}, scale=1.0]
% sommets du cube couleur
\coordinate (O) at (0,0,0); \coordinate (R) at (1,0,0); \coordinate (G) at (0,1,0); \coordinate (B) at (0,0,1);
\coordinate (RG) at (1,1,0); \coordinate (RB) at (1,0,1); \coordinate (GB) at (0,1,1); \coordinate (W) at (1,1,1);
\definecolor{cR}{rgb}{1,0,0}\definecolor{cG}{rgb}{0,1,0}\definecolor{cB}{rgb}{0,0,1}
\definecolor{cY}{rgb}{1,1,0}\definecolor{cM}{rgb}{1,0,1}\definecolor{cC}{rgb}{0,1,1}
\draw[gris, line width=0.6pt, dashed] (O) -- (G) -- (RG); \draw[gris, line width=0.6pt, dashed] (G) -- (GB);
\draw[navy, line width=0.8pt] (O) -- (R) -- (RG) -- (W) -- (GB) -- (B) -- cycle;
\draw[navy, line width=0.8pt] (R) -- (RB) -- (W); \draw[navy, line width=0.8pt] (B) -- (RB);
% axe des gris = direction de la luminance
\draw[line width=2.6pt, gris!70!navy] (O) -- (W);
\draw[fleche, navy, line width=1.6pt] (O) -- (0.62,0.62,0.62);
\node[font=\small, text=navy, anchor=west] (lum) at (1.45,1.05,0.62) {$Y$ : luminance};
\draw[navy, line width=0.5pt] (lum.west) -- (0.82,0.82,0.82);
\foreach \p/\c/\n/\pos in {O/black/noir/below left, R/cR/rouge/below, G/cG/vert/above left, B/cB/bleu/left, RG/cY/jaune/right, RB/cM/magenta/below right, GB/cC/cyan/above, W/white/blanc/above right}
  \filldraw[fill=\c, draw=navy, line width=0.6pt] (\p) circle (3.2pt) node[\pos=3pt, note] {\n};
\node[note, align=left, anchor=west] at (2.3,-0.2,0.2) {La diagonale noir $\to$ blanc porte la luminosité.\\ Les deux directions perpendiculaires\\ ($C_b$, $C_r$) ne portent que la teinte.};
\end{tikzpicture}
\end{document}
